\subsection{Execution details}
\label{app:execution}

For a native graph, the entrance pool concatenates node IDs from successive
upper layers, in increasing layer and ID order. A node on multiple upper
layers occurs repeatedly in the pool. A Java-compatible pseudorandom
generator seeded by the graph's node count samples that pool, rejecting
already accepted IDs until 64 distinct nodes (or all eligible nodes) have
been accepted. First acceptance determines table order. Imported graphs
retain their saved table. The order is independent of the query; no
recall or rescue label determines it.

Queues and answers use $(\hbox{distance},\hbox{node ID})$ order.
Admission compares $d<\tau$ before updating the local threshold;
local stopping compares $d>\tau$ between expansions.
An earlier admission may reach or exceed the tightened threshold, so
these tie rules differ. Initial stopping uses $\geq(1+\gamma)^2\tau_A$;
final continuation disables it.

Score validity, bottom-layer discovery, adjacency opening, and completion
are separate states. An upper-scored node can enter bottom traversal
without another evaluation. An already discovered neighbor is skipped
while the source cursor advances. Completion requires processing all
neighbors; stale queue references to completed nodes are then discarded.

F merges active, deferred, and current scout work into the main active
queue without replacing the main current node. Shared bottom cursors
preserve partial scans. F does not save an upper cursor if descent
exhausts its allowance: no bottom seed is installed, but computed scores
still contribute to the answer collector.

\subsection{Admission and stopping factorial}\label{app:factorial}
The six scout variants cross strict local admission (L) or all-neighbor
admission (A) with LOCAL, GLOBAL, or NONE distance stopping. Under LOCAL,
A and L have identical scored paths and answers on all 23,997 events;
the same agreement holds under GLOBAL. Under NONE, A returns 99.17698\%
mean recall, 24 severe events, and nine zero events, compared with
99.17031\%, 24, and ten for L. The unchanged severe count contains one
rescue and one harm. These retrospective results have descriptive
query-cluster intervals, not multiplicity-adjusted confirmation or new
timing. The original F policy is L\_LOCAL throughout the paper.

\subsection{Recall distribution and threshold checks}\label{app:thresholds}
\begin{table}[t]
\centering
\caption{Recall-tail sensitivity of C versus F on the 7,999-query cohort (23,997 graph--query events). Counts are descriptive; threshold 5 is the original severe endpoint.}
\label{tab:thresholds}
\small
\begin{tabular}{@{}rrrr@{}}
\toprule
Hits below & C & F & Rescue / harm \\
\midrule
1 & 25 & 10 & 17 / 2 \\
3 & 45 & 18 & 28 / 1 \\
5 & 65 & 26 & 41 / 2 \\
7 & 90 & 40 & 55 / 5 \\
9 & 280 & 218 & 95 / 33 \\
10 & 1542 & 1516 & 145 / 119 \\
\bottomrule
\end{tabular}
\end{table}

Table~\ref{tab:thresholds} recomputes saved hits from all 23,997 matched
pairs and checks the frozen totals. Counts are events, not independent
queries. Five remains the severe threshold; other cutoffs are descriptive,
not additional significance tests. The artifact retains the complete
zero-to-ten hit histograms and their source records.
